\documentclass{article}
\usepackage{amsmath,amssymb}
\usepackage{xurl}
\usepackage[hidelinks]{hyperref}
% Shared commands for notes in docs/paper-gaps/.

\DeclareUnicodeCharacter{2080}{\ifmmode{}_{0}\else\textsubscript{0}\fi}
\DeclareUnicodeCharacter{2081}{\ifmmode{}_{1}\else\textsubscript{1}\fi}
\DeclareUnicodeCharacter{2082}{\ifmmode{}_{2}\else\textsubscript{2}\fi}
\DeclareUnicodeCharacter{2099}{\ifmmode{}_{n}\else\textsubscript{n}\fi}

\newcommand{\Lean}{\textsc{Lean}}
\ExplSyntaxOn
\NewDocumentCommand{\leanid}{m}{
  \ifmmode
    \textsf{\detokenize{#1}}
  \else
    \group_begin:
    \urlstyle{sf}
    \tl_set:Nn \l_tmpa_tl {#1}
    \tl_replace_all:Nnn \l_tmpa_tl {\_} {_}
    \exp_args:NV \path \l_tmpa_tl
    \group_end:
  \fi
}
\ExplSyntaxOff
\providecommand{\tr}{\operatorname{tr}}
\newcommand{\tnleanIssueBase}{https://github.com/LionSR/TNLean/issues/}
\newcommand{\tnleanPrBase}{https://github.com/LionSR/TNLean/pull/}
\newcommand{\tnleanDocBase}{https://sirui-lu.com/TNLean/docs/}
\newcommand{\ghissue}[1]{\href{\tnleanIssueBase#1}{GitHub issue \##1}}
\newcommand{\ghpr}[1]{\href{\tnleanPrBase#1}{PR \##1}}
\newcommand{\doclink}[1]{\href{\tnleanDocBase#1.html}{\path{#1}}}

% Dirac notation and the identity operator, as in blueprint/src/macros/common.tex.
% \providecommand keeps note-local definitions authoritative.
\usepackage{dsfont}
\providecommand{\ket}[1]{|#1\rangle}
\providecommand{\bra}[1]{\langle #1|}
\providecommand{\braket}[2]{\langle #1 | #2 \rangle}
\providecommand{\ketbra}[2]{|#1\rangle\!\langle #2|}
\providecommand{\Id}{\mathds{1}}
\providecommand{\one}{\mathds{1}}
\providecommand{\C}{\mathbb{C}}
\providecommand{\MN}[1]{M_{#1}(\C)}
\providecommand{\MD}{M_D(\C)}

% External referencing for paper sources.  The build helper writes sanitized
% auxiliary files under build/paper-aux/ and excludes duplicated source labels.
\usepackage{xr}

% Import a generated paper auxiliary file with a caller-supplied label prefix.
% The candidate paths support builds from docs/paper-gaps/, the repository root,
% and build/paper-aux/ itself.
\newcommand{\importpaperaux}[2]{%
  \IfFileExists{../../build/paper-aux/#2.aux}{%
    \externaldocument[#1]{../../build/paper-aux/#2}%
  }{%
    \IfFileExists{build/paper-aux/#2.aux}{%
      \externaldocument[#1]{build/paper-aux/#2}%
    }{%
      \IfFileExists{#2.aux}{%
        \externaldocument[#1]{#2}%
      }{}%
    }%
  }%
}

% General external reference macros
\newcommand{\paperref}[2]{\ref{#1-#2}}
\newcommand{\papereqref}[2]{(\ref{#1-#2})}

% CPSV16 source (1606.00608 / MPDO-22-12-17-2.tex) external document and shortcuts
\importpaperaux{cpsv16-}{1606.00608/MPDO-22-12-17-2-tnlean}
\newcommand{\cpsvref}[1]{\paperref{cpsv16}{#1}}
\newcommand{\cpsveqref}[1]{\papereqref{cpsv16}{#1}}

% Verdict marker: \gapnote{<kind>}{<status>}. Kind is the policy
% classification (clarification, local-correction, scope-restriction,
% unfaithful, false-source, open-gap); status is open, wip (elimination
% underway), resolved, or historical. Machine-read by the site index;
% prints a verdict line.
\newcommand{\gapnote}[2]{\par\noindent\textbf{Verdict:} #1 (#2).\par}

\title{Full Support in MPU Normality and Theorem III.8}
\author{}
\date{2026-09-03}
\begin{document}
\maketitle
\gapnote{scope-restriction}{open}

\section{Source statement and surrounding convention}

The normal-tensor proposition of
Cirac--P\'erez-Garc\'ia--Schuch--Verstraete states that the normalized MPU
transfer matrix has exactly one nonzero eigenvalue, equal to one, and that the
normalized local tensor is normal
\cite[Proposition~III.1]{Cirac2017MPU}. The proposition does not state a
canonical-form or minimality hypothesis. The surrounding section has already
imposed the convention that, unless otherwise noted, the local tensor is first
replaced by one in canonical form
\cite[discussion preceding Proposition~III.1]{Cirac2017MPU}. The preceding
definition expresses the canonical tensor as a direct sum of irreducible
positive-dimensional blocks and permits its replacement by another
representative generating the same matrix product vectors.
% Source lines: paper_v2.tex 257--265, 319--326, and 344--354.

\section{The formal support issue}

The literal CPSV canonical-form data used in the formalization allow a
coisometric embedding of the retained direct sum into a larger ambient bond
space. This representation supports exact reconstruction but may contain an
ambient zero complement. The normality argument is valid for the intended
reduced representative obtained by removing that complement. The source does
not state this reduction as a separate theorem.

The formalization constructs the reduced representative in
\leanid{MPOTensor.IsMPU.exists_reduced_cfii_representative}. It has positive
bond dimension $D_{\mathrm{red}}\leq D$, admits full-support canonical-form-II
data, and generates the same periodic operators at every positive length. It
is a new representative; the theorem does not imply that the original ambient
tensor is normal.

The constructed representative is written in the ambient coordinates of its
unique retained block. In those coordinates the ambient right fixed matrix is
the block's own diagonal positive fixed matrix, normalized to trace one. This
is the matrix denoted $\rho$ in the source canonical form II, namely the
diagonal positive definite right fixed matrix of the normalized transfer map
with $\tr(\rho)=1$; it is exhibited in ambient coordinates in
Eq.~\ref{eq:mpu_canonical_form_full_support_ambient_fixed_point} below. The
construction
therefore inhabits the bundled canonical-form-II presentation
\leanid{MPOTensor.IsMPUCanonicalFormII}. This inhabitance is specific to the
constructed gauge: canonical-form-II data together with full support do not by
themselves make the ambient right fixed matrix diagonal, because a
non-monomial unitary change of ambient coordinates carries a non-scalar
diagonal fixed matrix to a nondiagonal one. No such general projection is
claimed.

This distinction is necessary. Given any MPU representative, adjoin a zero
virtual direct summand. The periodic contractions do not change, so the
enlarged tensor generates the same unitary operators at every length. However,
the projection onto the original virtual summand is a nontrivial invariant
projection. The enlarged ambient tensor is therefore reducible and hence not
normal. Thus the bare assertion ``MPU implies normal tensor'' is false for
nonminimal ambient representatives.

Let $D$ be the ambient bond dimension, and let $D_k$ be the dimensions of the
retained canonical blocks. Full support means
\begin{align}
    \sum_k D_k & = D.
    \label{eq:mpu_canonical_form_full_support_total_dimension}
\end{align}
If there is exactly one block,
Eq.~\ref{eq:mpu_canonical_form_full_support_total_dimension} states that this
block has dimension $D$. Positivity of the retained block dimensions and the
total-dimension bound exclude every additional retained block. The ambient
coisometry is then square and therefore unitary. Block assembly followed by
unitary conjugation transports irreducibility from the normal block to the
ambient tensor.

The normalized-flattening construction is formalized by
\leanid{MPOTensor.IsMPU.exists_reduced_normalizedFlattening_cfii}. Starting
from the nonzero irreducible reduction, it transfers the shifted trace
identities to the retained direct sum. Every transfer eigenvalue of a retained
block embeds into the ambient transfer spectrum. Its positive Perron value and
each peripheral value therefore equal one, so every retained block is normal.
The canonical-form-II gauge preserves their total bond dimension. A final
gauge transformation by the inclusion of the unique retained block places the
representative in the coordinates carrying the diagonal ambient right fixed
matrix.

\section{The source factors and Theorem III.8}

Let $\mathcal U$ be an MPU tensor with physical dimension $d$ and bond
dimension $D$, and let $E$ be the transfer matrix of the normalized tensor
$d^{-1/2}\mathcal U$. After the normal-tensor proposition, the source assumes
without loss of generality that this normalized tensor is in canonical form II
\cite[discussion following Proposition~III.1]{Cirac2017MPU}. Fix such data with
full support.
% Source lines: paper_v2.tex 350--356.

The reduced-representative theorem justifies this replacement for the
positive-length periodic operator family. It does not identify the compact-SVD
cuts or selected source factors of the new tensor with those of the original
auxiliary presentation. A result that retains the source cuts of a named tensor
therefore still needs source data for that tensor or a separate comparison
theorem.

Let $\mathcal A_1$ be the unique block, let $V$ include that block into the
ambient bond space, and let $\Lambda$ be its diagonal positive right fixed
point normalized to have trace one. Define
\begin{align}
    \rho & = V\Lambda V^\dagger.
    \label{eq:mpu_canonical_form_full_support_ambient_fixed_point}
\end{align}

The source factors use $\rho$. In the first singular value decomposition, let
$V_1$ be the coisometry in $M_1=V_1^\dagger D_1U_1$. The first source cut has
row order $(i,\beta)$, namely upper physical followed by right auxiliary. Its
weight and compressed positive metric are therefore
\begin{align}
    W_\rho & = \mathbf{1}_d\otimes\rho, \\
    G      & = V_1W_\rho V_1^\dagger.
    \label{eq:mpu_canonical_form_full_support_source_metric_physical_virtual}
\end{align}
The normalization factor in $X_1$ is $G^{-1/2}$, not $G$.

Existence of the inverse square root requires only that
$V_1W_\rho V_1^\dagger$ be positive definite. It does not force $\rho$ to be
positive definite on the full ambient bond space. The current TNLean
source-factor construction assumes ambient positive definiteness, and full
support supplies this stronger condition
\cite[Section~III, source-factor construction]{Cirac2017MPU}. The same factors
define the operators $u$ and $v$ in Lemma~III.7 and Theorem~III.8 of
Ref.~\cite{Cirac2017MPU}.
% Source lines: paper_v2.tex 479--505 and 532--601.

Ambient positive definiteness is therefore a restriction of the current formal
construction. It is not a consequence of invertibility of the compressed
source metric or of the bare MPU predicate.

Let $(\Phi\rvert)=\operatorname{vec}(\mathbf{1}_D)^{\mathsf T}$ be the
vectorized identity bra. The supported formal route has two bond-dimension
branches. If $D=1$, the theorem
\leanid{MPOTensor.IsMPU.normalized_transfer_matrix_eq_one_fin_one} gives
\begin{align}
    J    & = 1,
    \label{eq:mpu_canonical_form_full_support_fin_one_exponent}    \\
    \rho & = \mathbf{1}_1,
    \label{eq:mpu_canonical_form_full_support_fin_one_fixed_point} \\
    E^J  & = \lvert\rho)(\Phi\rvert.
    \label{eq:mpu_canonical_form_full_support_fin_one_stabilization}
\end{align}
If $D>1$, the route assumes chosen canonical-form-II data $\mathcal C$ such
that
\begin{align}
    \mathcal C
     & \in\operatorname{CFII}(d^{-1/2}\mathcal U),
    \label{eq:mpu_canonical_form_full_support_cfii_membership} \\
    \sum_k D_k
     & = D.
    \label{eq:mpu_canonical_form_full_support_cfii_full_support}
\end{align}
The second identity is the full-support condition. The theorem
\leanid{MPOTensor.IsMPU.normalized_transfer_power_eq_vecMulVec_of_reduced_cfii}
then supplies the unique block $\mathcal A_1$, the matrices $\Lambda$ and
$\rho$, and the identities
\begin{align}
    \rho
     & = V\Lambda V^\dagger,
    \label{eq:mpu_canonical_form_full_support_reduced_fixed_point}    \\
    \rho
     & >0,
    \label{eq:mpu_canonical_form_full_support_fixed_point_positive}   \\
    \tr(\rho)
     & =1,
    \label{eq:mpu_canonical_form_full_support_fixed_point_trace}      \\
    E\lvert\rho)
     & =\lvert\rho),
    \label{eq:mpu_canonical_form_full_support_right_fixed_vector}     \\
    (\Phi\rvert E
     & =(\Phi\rvert,
    \label{eq:mpu_canonical_form_full_support_left_fixed_vector}      \\
    E^J
     & =\lvert\rho)(\Phi\rvert,
    \label{eq:mpu_canonical_form_full_support_transfer_stabilization} \\
    J
     & =D^2-1.
    \label{eq:mpu_canonical_form_full_support_stabilization_exponent}
\end{align}
In both branches, the positive trace-one matrix $\rho$ is the source weight
used to construct the factors and the operators $u$ and $v$.

The explicit reduced-CFII and full-support hypotheses resolve the
ambient-positivity question tracked by \ghissue{5855}. For a supplied fixed
pair whose matrix $\rho$ is diagonal, the arbitrary-$K$ complete source-$u$
network is formalized by
\leanid{MPOTensor.SourceFactors.sourceU_gram_eq_normalized_input_tail}.
Input unitarity gives the isometry and rank bound in
\leanid{MPOTensor.SourceFactors.sourceU_isIsometry_of_isMPU} and
\leanid{MPOTensor.SourceFactors.mul_self_le_rightRank_mul_leftRank_of_isMPU};
\leanid{MPOTensor.IsMPUCanonicalFormII.sourceU_isIsometry} applies this result
to the compact-SVD source factors of a tensor in canonical form II, where the
weight and the transfer power come from the convention. The raw trace-one transfer-power specialization is
\leanid{MPOTensor.SourceFactors.sourceU_gram_eq_normalized_mpo_input_tail_of_isDiag}.
These proofs keep the source factors built from $\rho$, use
$\rho^{\mathsf T}=\rho$, and retain one $K$-site interior between two
unblocked endpoints. This is the supplied-fixed-pair form of Lemma~III.7 of
Ref.~\cite{Cirac2017MPU}; it does not construct such source data for an
arbitrary MPU representative.

The complete source-$v$ equivalence from the proof of Theorem~III.8 of
Ref.~\cite{Cirac2017MPU} is formalized as
\leanid{MPOTensor.sourceVDressedGram_iff_simple2_transpose_fixed_pair}; it
needs no fixed-pair datum beyond a positive definite weight, and it carries
the inserted-vector transpose recorded in
\texttt{mpu\_source\_cut\_orientation.tex}. The source's diagonal convention
then gives \leanid{MPOTensor.IsMPUCanonicalFormII.sourceV_isIsometry} for the
gate built from that same weight.
% Source lines: paper_v2.tex 545--557 and 577--600.
% Source equation labels: uUnitary and vUnitary.

For $D>1$, the admissible source results assume reduced canonical-form-II data
with full support for the same tensor $\mathcal U$ whose cuts, ranks, and source
operators appear in their statements. For $D=1$, the one-dimensional
stabilization theorem supplies the source pair directly. In either branch, the
datum supplies $J>0$, a positive trace-one matrix $\rho$, and
\begin{align}
    E^J & = \lvert\rho)(\Phi\rvert.
    \label{eq:mpu_canonical_form_full_support_supplied_stabilization}
\end{align}
The source-$u$ theorem additionally assumes that $\rho$ is diagonal, as in the
source's canonical-form-II coordinates. Supplying these same-tensor diagonal
source data is part of the construction tracked by \ghissue{7012}; the later
unrestricted representative remains tracked by \ghissue{7127}. The admissible
source-isometry and simple-tensor equivalence nodes cite Lemma~III.7 and
Theorem~III.8 of Ref.~\cite{Cirac2017MPU} directly. Both retain the source
factors of $\mathcal U$.
% Source lines: paper_v2.tex 545--601.

Blocking occurs only within the supplied-fixed-pair source-$u$ calculation.
Let $J$ be the supplied stabilization exponent and define $K=JD^2$. The
reflected transfer identity converts the power at $J$ into the direct
$K$-site double-layer fixed pair. The calculation then keeps this one interior
between the two unblocked endpoint letters. It does not replace the source
cuts, ranks, or operators of $\mathcal U$, and it does not transport source
data across blocking.

\section{Downstream standard-form, index, and symmetry scope}

Chapter~28 of the blueprint records the preceding assumptions as a reduced
full-support source datum for the same tensor. This datum is an explicit input,
not an existence theorem. The labels \texttt{lemuisometry} and
\texttt{ThmFund1} are no longer among the unformalized source entries: both
are now stated and proved for a tensor satisfying the bundled
canonical-form-II presentation, which is the standing hypothesis the source
fixes once before Lemma~III.7 and never repeats, and both carry the source's
hypothesis set. The chapter retains unrestricted, unformalized source entries
under the original labels \texttt{SF}, \texttt{def:index},
\texttt{index-well-defined}, \texttt{lem:mpu\_index\_tensor}, and
\texttt{lem:mpu\_index\_composition}. The label
\texttt{lem:mpu\_admissible\_source\_u\_isometry} likewise no longer
records a restriction: it is the reusable supplied-fixed-pair step, stated for
an arbitrary positive diagonal weight and an arbitrary stabilization exponent,
through which \texttt{lemuisometry} is proved. The remaining restricted
full-support statements include the distinct labels
\texttt{def:mpu\_admissible\_index},
\texttt{prop:mpu\_admissible\_index\_well\_defined},
\texttt{lem:mpu\_admissible\_index\_tensor}, and
\texttt{lem:mpu\_admissible\_index\_composition}.

Chapter~28 now contains thirty-six labelled \texttt{mpu\_admissible} entries.
Thirty-five are restricted forms of results or definitions in
Ref.~\cite{Cirac2017MPU}; each has a direct citation or an exact source-line
locator. The remaining entry,
\texttt{def:mpu\_admissible\_equivalence\_datum}, is an auxiliary hypothesis
introduced in the blueprint. It cites the source path arguments motivating its
conditions but does not attribute the datum itself to the paper.

The unrestricted source Fundamental Theorem of MPU remains under
\texttt{FundamentalMPU} and depends only on the source standard form
\texttt{SF}. The former admissible duplicates of these statements and of
\texttt{ThmFund1} have been removed: the same-tensor converse supplies a
canonical-form-II presentation recording each supplied fixed matrix.
The unrestricted source Index Theorem remains under
\texttt{IndexTh}. It depends on the source index, source well-definedness,
the unrestricted source-cut formulations of tensor-product and composition
additivity, source continuity, and
source equivalence.

The earlier admissible Index Theorem still records operation-specific source
data for its other conclusions. Its independent-product clause remains
restricted. The composition-additivity entry, however, now states the proved
representative-index theorem with no six-block source-data hypotheses. A
blocking length is admissible for statements naming source factors only if
the actual blocked tensor is simple and carries its own source datum.
For a tensor already in canonical form II, the two index entries now use
\leanid{MPOTensor.IsMPUCanonicalFormII.index}. Positive blocking preserves
that presentation, and simplicity gives $r_k\ell_k=d^{2k}$. Thus index
well-definedness no longer requires separately supplied source data at each
simple block. The value is independent of the positive simple blocking and
of the recorded canonical-form-II presentation of the same tensor; it is
preserved by positive blocking and changes sign under physical adjunction.
The representative index is now defined for every positive-dimensional MPU
through a reduced canonical-form-II presentation. The unitary virtual-gauge
theorem identifies the right and left source-cut ranks of canonical-form-II
presentations with equal periodic families, even across initially different
ambient bond dimensions. A common positive simple blocking then identifies
their indices. Consequently the representative index is independent of the
chosen reduction and of the ambient bond dimension
(\leanid{MPOTensor.IsMPUCanonicalFormII.sourceRanks_eq_of_mpo_eq},
\leanid{MPOTensor.IsMPUCanonicalFormII.index_eq_of_mpo_eq},
\leanid{MPOTensor.IsMPU.index_eq_canonical_representative},
\leanid{MPOTensor.IsMPU.index_eq_of_mpo_eq}). This does not assert that
rectangular reduction preserves the two source-cut ranks of the unreduced
tensor itself, or identify this representative definition with the printed
unrestricted source-cut formula without further argument.
For the representative index, composition additivity now holds for arbitrary
positive-dimensional MPUs (\leanid{MPOTensor.IsMPU.index_mulTensor}). The
canonical-form-II proof chooses a common simple blocking rather than assuming
one; this public theorem does not identify the raw source cuts of an unreduced
tensor with those of its representative. Independent tensor-product
additivity is a separate development and is not part of this composition
argument.

The published composition proof first uses $\mathcal U_k$,
$\mathcal U'_k$, and $\mathcal W_k$, then places $\mathcal U_{2k}$ and
$\mathcal U'_{2k}$ in standard form and computes the ranks of
$\mathcal W_{4k}$ \cite{Cirac2017MPU}. The earlier admissible formulation
supplied source data for all six exact blocked tensors. The composition entry
has now been replaced by the representative-index theorem, whose proof does
not use those six assumptions. It reduces the product tensor to a
canonical-form-II representative with the same positive-length periodic
operators (\leanid{MPOTensor.exists_isReduction_of_mpo_eq_cfii}). Rectangular
reduction only decreases the two source ranks. At a common simple blocking,
the four-site composition upper bounds and the exact rank product of the
reduced representative force both bounds to be equalities
(\leanid{MPOTensor.IsMPUCanonicalFormII.compositionRanks_of_isReduction}).
This proves canonical-representative composition additivity
(\leanid{MPOTensor.IsMPUCanonicalFormII.index_eq_add_of_mpo_eq_mulTensor}),
which descends to the public representative index.
The supplied-fixed-pair source-$u$
contraction and the complete source-$v$ network are the steps through which
\texttt{lemuisometry} and \texttt{ThmFund1} are proved for a tensor in
canonical form II. Together with preservation under positive blocking, they
establish the index of a canonical-form-II representative described above.
Statements that name selected source factors or raw source cuts of a
transformed tensor retain their separate construction obligations.
% Source lines: paper_v2.tex 824--865.

The same separation between source statements and scoped statements extends
through the continuity and symmetry subtree. The continuity labels
\texttt{prop:continuity-index} and \texttt{coruvSF} have the scoped
counterparts \texttt{prop:mpu\_admissible\_continuity\_index} and
\texttt{cor:mpu\_admissible\_continuous\_standard\_form}. The latter is the
restricted-source form of the continuous standard-form corollary
\cite[lines~867--910]{Cirac2017MPU}.

For any tensor $\mathcal A$, let $\mathcal A^\sharp$ denote its
physical-adjoint comparison tensor. If $\mathcal A_k$ is the actual $k$-site
block, then $\mathcal A_k^\sharp$ is the physical adjoint of that blocked
tensor.

For local symmetry and conjugation classification, the source labels
\texttt{prop1conj}, \texttt{corconj},
\texttt{cor:sympl-to-orthogonal}, \texttt{PropChiral}, and
\texttt{prop:local-char-transposition} have the scoped counterparts
\texttt{prop:mpu\_admissible\_conjugation\_symmetry},
\texttt{prop:mpu\_admissible\_conjugation\_standard\_forms},
\texttt{cor:mpu\_admissible\_symplectic\_to\_orthogonal},
\texttt{prop:mpu\_admissible\_local\_time\_reversal}, and
\texttt{prop:mpu\_admissible\_local\_transposition}. The unrestricted local
characterizations reproduce the propositions of Ref.~\cite{Cirac2017MPU} and
therefore do not add an index-zero hypothesis; the index obstruction belongs
to the surrounding source discussion. Their admissible counterparts also do
not assume index zero.

For the supplied two-site tensor and its physical-adjoint or transposed
comparison, symmetry identifies the two generated MPUs. The Admissible
Fundamental Theorem then identifies their source spaces. The relevant operation
exchanges the source cuts, giving
\begin{align}
    r(\mathcal U_2)
     & =r(\mathcal U_2^\sharp),
    \label{eq:mpu_canonical_form_full_support_adjoint_right_rank} \\
    r(\mathcal U_2^\sharp)
     & =\ell(\mathcal U_2).
    \label{eq:mpu_canonical_form_full_support_adjoint_cut_exchange}
\end{align}
The transposed comparison tensor satisfies the analogous identities. Hence
$\operatorname{ind}(\mathcal U_2)=0$ follows directly from the source ranks,
and admissible blocking independence yields
$\operatorname{ind}(\mathcal U)=0$. The supplied source and comparison data
therefore imply the obstruction; it is not an external hypothesis. This
argument follows the discussion preceding Propositions~IV.5 and~IV.6 of
Ref.~\cite{Cirac2017MPU}.
% Source lines: paper_v2.tex 1201--1209 and 1257--1261.

The source definitions \texttt{def:strictly-equivalent-symmetry} and
\texttt{def:equivalent-symmetry}, and the source criterion
\texttt{lem:equivalent-symmetry}, have the scoped labels
\texttt{def:mpu\_admissible\_strict\_symmetry\_equivalence},
\texttt{def:mpu\_admissible\_symmetry\_equivalence}, and
\texttt{lem:mpu\_admissible\_symmetry\_path\_criterion}. The source equation
at paper lines~1345--1353 repeats the symmetry-transformed operator on both
sides. The blueprint instead retains the intended invariance condition
\begin{align}
    \mathcal S[W^{(N)}(p)]
     & =W^{(N)}(p).
    \label{eq:mpu_canonical_form_full_support_symmetry_invariance}
\end{align}
% **Local fix:** corrected the repeated right-hand side in the source display.

The conjugation-phase results \texttt{conjphs},
\texttt{cor:all-equivalent-under-conjugation},
\texttt{prop:conjsym-strict-equiv}, and
\texttt{cor:equiv-under-conjugation} have the scoped labels
\texttt{thm:mpu\_admissible\_conjugation\_odd\_rank},
\texttt{cor:mpu\_admissible\_all\_conjugation\_equivalent},
\texttt{thm:mpu\_admissible\_conjugation\_even\_rank}, and
\texttt{cor:mpu\_admissible\_equiv\_under\_conjugation}. The unrestricted
all-equivalent corollary retains the source condition that either $r$ or
$\ell$ is odd.

The admissible odd-rank determinant criterion assumes reduced full-support
source data for each original one-site endpoint and its conjugate comparison,
in addition to the blocked path datum. These assumptions define the endpoint
standard-form unitaries $u_i$ and $v_i$. The final admissible equivalence
corollary retains both source branches. If one original source rank is odd,
only Case~I occurs, but this does not determine the parity of the source ranks
after blocking.

The identities $r^{(2)}=dr$ and $\ell^{(2)}=d\ell$ use admissible blocking
independence and therefore require source data for the simple one-site and
two-site endpoints. Both conjugate two-site endpoint tensors also carry source
data. If a blocked rank is odd, only the blocked product determinant is needed,
and the odd-rank criterion applies. If both blocked ranks are even, including
the case in which an even physical dimension multiplies an odd source rank,
blocking preserves Case~I. Parity then makes the two separate blocked
determinants equal to one, so the even-rank criterion applies. If both original
ranks are even, the blocked ranks remain even, and the even-rank criterion
distinguishes Case~I from Case~II.

The admissible odd-original-rank all-equivalent corollary uses this complete
classification rather than applying the odd-rank theorem after blocking
without a parity check. Its hypotheses supply source data for the one-site and
two-site endpoints and both conjugate comparisons. The admissible
$U_1$--$U_2$ conjugation classification repeats these four data requirements
for each of $U_1$, $U_2$, $\widetilde U_1$, and $\widetilde U_2$. A path datum
at an unrelated length $m$ does not replace the endpoint data.
% **Local fix:** paper_v2.tex 1554--1559 writes the two separate determinant
% identities before the parity distinction. The admissible proof uses them
% only in the both-blocked-even branch, where the exponent is even.
% Source lines: paper_v2.tex 1385--1460 and 1548--1567.

For the time-reversal sign, the source labels
\texttt{eq:sigma-from-trace},
\texttt{lem:mpu\_sigma\_gauge\_invariant},
\texttt{prop:sigma-well-defined},
\texttt{prop:sigma1-and-sigma2-equal}, and
\texttt{cor:sigma-equivalence-necessary} have the scoped labels
\texttt{lem:mpu\_admissible\_sigma\_trace},
\texttt{lem:mpu\_admissible\_sigma\_gauge\_invariant},
\texttt{prop:mpu\_admissible\_sigma\_well\_defined},
\texttt{prop:mpu\_admissible\_sigma\_blocking\_parity}, and
\texttt{cor:mpu\_admissible\_sigma\_equivalence\_necessary}. The two MPS
comparison results
\texttt{prop:mps-class-eq-mpu-class-for-conjugation} and
\texttt{prop:dagger-mps-class-and-mpu-class-equal} have the scoped labels
\texttt{prop:mpu\_admissible\_mps\_conjugation\_class} and
\texttt{prop:mpu\_admissible\_mps\_time\_reversal\_class}.

The source example and swap labels
\texttt{ex:all-phases-under-conjugation}, \texttt{example:czx},
\texttt{ex:dagger-all-but-czx}, \texttt{lemma:sym-trafo-swap},
\texttt{prop:U2-U3-trivial-conjugation},
\texttt{prop:U2-U3-trivial-tr-transpose},
\texttt{prop:eqiv-U1-U2-conjugate}, and
\texttt{prop:U1-U2-equiv-ancillatrick} have the scoped counterparts
\texttt{ex:mpu\_admissible\_all\_conjugation\_phases},
\texttt{ex:mpu\_admissible\_czx},
\texttt{ex:mpu\_admissible\_time\_reversal\_classes},
\texttt{lem:mpu\_admissible\_swap\_symmetry\_equivalence},
\texttt{prop:mpu\_admissible\_u2\_u3\_conjugation},
\texttt{prop:mpu\_admissible\_u2\_u3\_time\_reversal\_transposition},
\texttt{prop:mpu\_admissible\_u1\_u2\_conjugation}, and
\texttt{prop:mpu\_admissible\_u1\_u2\_ancilla}. The unrestricted CZX entry
defines the Hadamard matrix used in its tensor formula. Comment-only source
locators identify the exact passages for the unitary normal-form lemma, the
three topological-insulator-inspired MPUs, and their swap transformation.

Continuity needs more than endpoint data. Chapter~28 therefore introduces an
admissible equivalence and path datum. It contains a chosen strict-equivalence
path, one common positive blocking length $m$, and continuously varying reduced
full-support source data for the actual blocked path tensor
$\mathcal W(p)_m$ at every point. The tensor $\mathcal W(p)_m$ must itself be
simple; it cannot be replaced by an equivalent canonical or reduced
representative. The actual blocked endpoints have their own data. The same
condition applies after adjoining ancillas and blocking.

This datum is a hypothesis on a chosen witness, not an existence statement. It
does not automatically include the source bound $m\leq D^4$. The continuous
standard-form corollary assumes that bound. For a datum of length $m$, the
standard-form families describe the actual blocked endpoint MPUs
$U_1^{(mN)}$ and $U_2^{(mN)}$. They describe the unblocked families
$U_1^{(N)}$ and $U_2^{(N)}$ only when $m=1$. Every admissible strict-phase
theorem about the original endpoints therefore requires a path datum of
blocking length one. General admissible equivalence may still be witnessed by
strict equivalence of supplied blocked endpoints.

For a path preserving a specified symmetry $\mathcal S$, admissibility also
includes continuously varying source data for the fixed family of conjugate,
physical-adjoint, transposed, and $\mathcal S$-image comparison tensors of the
actual blocked path. Every member of this family is simple and has an
independent datum, whether or not a later proof uses that comparison. No
symmetry operation is assumed to preserve source data automatically.

Under this convention, admissible equivalence implies equality of the index.
The converse is asserted only when the standard-form interpolation has an
admissible path datum. The admissible standard-form path criterion at length
$m$ concerns the blocked endpoint families $U_1^{(mN)}$ and $U_2^{(mN)}$. Its
strict-phase descendants concern the original endpoint families only when
$m=1$. Thus an ordinary strict-equivalence or symmetry-equivalence path is
never silently promoted to an admissible path, and a blocked standard form is
not treated as a standard form of an unblocked endpoint.

For a time-reversal-invariant tensor $\mathcal U$, let
$\sigma^{(k)}\in\{\pm1\}$ be the sign in the local time-reversal relation for a
standard form of the actual $k$-site block $\mathcal U_k$.

The source proofs of the time-reversal and transposition characterizations
apply the fundamental theorem after blocking two sites
\cite{Cirac2017MPU}. Each corresponding admissible statement therefore
supplies source data for the actual one-site endpoint tensor, its actual
two-site block, and the physical-adjoint or transposed comparison tensor of
that block. Defining both $\sigma^{(1)}$ and $\sigma^{(2)}$ requires source
data for $\mathcal U$, $\mathcal U_2$, $\mathcal U_2^\sharp$,
$\mathcal U_4$, and $\mathcal U_4^\sharp$.

Gauge invariance of the sign is separate and follows directly from the trace
formula, without a path hypothesis. The blocking-parity proposition uses this
gauge-invariance lemma and the explicit unwinding argument of the source; it
does not use pathwise constancy.

To compare both signs along an unblocked path $\mathcal W(p)$, pathwise
constancy is applied separately to the continuous families
$\mathcal W(p)$, $\mathcal W(p)_2$, $\mathcal W(p)_4$,
$\mathcal W(p)_2^\sharp$, and $\mathcal W(p)_4^\sharp$. This is a length-one
admissible path. After adjoining ancillas, the same five continuous families
are required for the actual ancilla-enlarged blocked witness path, considered
as a length-one path in its own right rather than only through its endpoints.

The resulting non-strict conclusion compares only $\sigma^{(2)}$ of the
supplied witness endpoints $\mathcal W_{\mathrm{anc}}(0)$ and
$\mathcal W_{\mathrm{anc}}(1)$. It does not compare $\sigma^{(2)}$ of the
original MPU tensors. Promoting the witness-level equality to an invariant of
the original tensors requires preservation of the selected source data under
adjoining identity ancillas. The closed \ghissue{6138} proves preservation of
CFII and full support but does not identify the selected compact-SVD factors.
That remaining boundary is tracked by \ghissue{6136}.

All later time-reversal descendants and examples repeat the necessary endpoint
and comparison data after blocking or adjoining ancillas. The swap
transformation and its two dependent strict-phase statements require the
actual source-constructed interpolation to have an admissible symmetry-path
datum of blocking length one. This datum includes the path and all four fixed
comparison paths: conjugate, physical-adjoint, transposed, and symmetry-image.
No preservation under these operations is asserted.

In the admissible $U_1$--$U_2$ ancilla statement, the supplied length-one path
begins at named ancilla-enlarged blocked endpoints
$\widehat{\mathcal U}_1$ and $\widehat{\mathcal U}_2$. Only these endpoints are
claimed to be admissibly strictly equivalent. The statement does not assert
equivalence of the original ancilla-free tensors $\widetilde U_1$ and
$\widetilde U_2$. The construction adjoins equal ancillas of dimension $d$,
which need not satisfy the coprimality condition in the blueprint definition
of equivalence when $d>1$.

For bond dimension one, the normalized transfer matrix supplies the source pair
directly. For every positive bond dimension, the reduced-representative theorem
now supplies a full-support canonical-form-II tensor generating the same
positive-length periodic operators. It does not identify the source data of the
original tensor with those of the reduced representative. It also does not
prove preservation of source data under tensor products, composition,
conjugation, adjoint, transposition, or continuous variation. The public
representative index nevertheless satisfies composition additivity, whose
proof does not require such preservation.

Consequently, the remaining admissible Chapter~28 statements retain their
explicitly supplied tensors, comparison tensors, and paths. Parallel
source-labelled statements retain the hypotheses of Ref.~\cite{Cirac2017MPU};
those not yet proved remain unformalized. No restricted statement is presented
under a source label.

\section{Elimination plan}

The native source-data scope tracker is \ghissue{6136}. The
dimension-bounded full-support representative is constructed by
\ghissue{7065} and the two theorems identified above. This closes
representative existence but does not transport selected compact-SVD factors
between unequal bond spaces. The supplied-fixed-pair source-$u$ network from
\ghissue{5982} and the source-$v$ conclusion are formalized. The admissible
source-unitary and standard-form construction remains tracked by
\ghissue{7012}, while \ghissue{7127} owns the later unrestricted standard-form
representative.

Issue~\ghissue{6139} treats tensor products, composition, and conjugate,
adjoint, transposed, and other transformed comparison tensors. It removes the
additional endpoint and comparison-tensor assumptions from the admissible
local-symmetry and example nodes. Issue~\ghissue{6140} constructs continuous
source data along strict-equivalence paths and their four fixed
symmetry-transformed comparison paths. It removes the path hypotheses from the
admissible continuity, symmetry-equivalence, phase-classification, swap, and
MPS-comparison nodes.

The twenty-nine source/scoped pairs in the continuity and symmetry subtree are
merged only after the relevant branches close. Until then, the unproved source
labels remain unrestricted statements, and the corresponding reduced-source or
path-dependent statements remain under their \texttt{mpu\_admissible} labels.
The supplied-fixed-pair contraction does not by itself construct the source
data required to remove these restrictions.

\subsection*{The pointwise family is one connected group}

The supplied fixed-pair converse separates restrictions on the same tensor
from those on derived tensors and continuous paths. Of the original thirty-nine
\path{mpu_admissible} entries of Chapter~28, twenty-three are
\path{def:mpu_admissible_equivalence_datum} together with the entries
whose dependencies cite it. These carry continuously varying path data and
belong to \ghissue{6140}. The original sixteen pointwise entries are now
thirteen after the three same-tensor duplicates were removed.

One of the thirteen,
\texttt{lem:mpu\_admissible\_source\_u\_isometry}, is not a duplicate. It is
the generic step, stated for an arbitrary positive diagonal weight $\rho$ and
an arbitrary exponent $K$ with $E^{K}=\lvert\rho)(\Phi\rvert$, and its formal
statements still carry both inputs. The source-labelled Lemma~III.7 entry is
proved from it together with the stabilized power of the bundled presentation,
so it remains as the reusable step.

The three removed entries ---
\path{thm:mpu_admissible_simple_tensor_equivalence},
\path{def:mpu_admissible_standard_form}, and
\path{thm:mpu_admissible_fundamental} --- restricted only the tensors that
their source-labelled counterparts already place under the standing
convention. The converse supplies a presentation on the same tensor with
exactly the fixed matrix used to construct its source operators. Their
retained dependents now cite the converse and the source-labelled statements.

The remaining twelve restrict a tensor other than the one the source entry is
about: an exact physical block $\mathcal U_k$, a physical-adjoint, transposed,
or conjugate comparison tensor, a tensor product, or a composition. Assuming
the presentation there is not the presentation on $\mathcal U$; it is the
assertion that the presentation survives those operations. For blocking the
source asserts exactly that, writing that the blocked tensor is again normal
and in canonical form II \cite[line~356]{Cirac2017MPU}. Positive blocking now
preserves the bundled canonical-form-II presentation, including the same
ambient fixed matrix:
\leanid{MPOTensor.IsMPUCanonicalFormII.blockTensor}. Physical adjunction
and independent tensor products preserve the bundled presentation by
\leanid{MPOTensor.IsMPUCanonicalFormII.physicalAdjointTensor} and
\leanid{MPOTensor.IsMPUCanonicalFormII.tensorProduct}, respectively. Their
downstream derived-tensor entries have not yet been consolidated. Separate
conjugation and transposition comparisons, symmetry transforms, and
composition source-factor statements retain distinct obligations under
\ghissue{6139}. Representative-index composition additivity is treated by
rectangular reduction and rank saturation instead.

Before the supplied-pair converse and the preservation constructors, the
recorded census contained sixteen declarations mentioning the bundled
presentation: the structure itself, thirteen consumers, and two constructors.
The thirteen consumers establish, for the same tensor, ambient full support,
nonempty bond and physical spaces, normality of the normalized flattening,
forced block contractions, the simplicity equivalence, the two stabilized
transfer identities, the normalized diagonal power, and the four source-factor
isometry and rank statements. The two constructions produce
a presentation from canonical data for a tensor equal to a normalized
flattening, and from an arbitrary matrix product unitary by dimension
reduction. None of those sixteen carried a presentation on $\mathcal U$ to a
presentation on a physical block, an adjoint, a transpose, a conjugate, a
tensor product, or a composition, which is what the twelve entries would need;
for the composition entry the transport has to run through a reduced
representative, as the last part of this subsection shows. Subsequent
constructors supply preservation of the bundled presentation under positive
blocking, physical adjunction, and independent tensor products. The converse
constructs a presentation on the tensor carrying the supplied fixed pair.
None of these constructions identifies the source cuts or chosen source
factors of a derived tensor with those of $\mathcal U$.

The three same-tensor entries have been consolidated with their source-labelled
counterparts under \ghissue{7660}. Eleven retained entries had cited them:
seven pointwise entries, and the four path entries
\path{prop:mpu_admissible_continuity_index},
\path{thm:mpu_admissible_index},
\path{cor:mpu_admissible_continuous_standard_form}, and
\path{lem:mpu_admissible_symmetry_path_criterion}. The tensors of all
eleven carry a supplied fixed pair. Their statements retain those pairs and
all operation-specific and pathwise hypotheses. The redirected dependencies
use the implication
\begin{align}
    \text{unitary periodic family},\ E^{J}=\lvert\rho)(\Phi\rvert
     & \implies\text{canonical form II}
    \label{eq:mpu_canonical_form_full_support_datum_to_convention}
\end{align}
is now recorded by
\leanid{MPOTensor.IsMPU.exists_canonicalFormII_of_supplied_transfer_power}.
The resulting presentation is on the same tensor and records the supplied
$\rho$ itself, so the source operators formed from $\rho$ are unchanged. This
does not introduce a parallel definition of canonical form.

There are two independent mathematical stages. The first, concerning the four
original entries whose restricted tensor is $\mathcal U$ itself, is proved
below. Its generic source-$u$ lemma is retained, and the three duplicate
entries have been removed. The second concerns the twelve entries whose
restricted tensor is derived from $\mathcal U$. Positive
blocking, physical adjunction, and independent tensor products now preserve
the presentation, but their downstream uses have not been consolidated; the
other operations remain open.

\begin{itemize}
    \item The proved converse of the stabilization theorem, namely
          (\ref{eq:mpu_canonical_form_full_support_datum_to_convention}): for
          a tensor whose periodic operators are all unitary, a positive
          diagonal trace-one $\rho$ with $E^{J}=\lvert\rho)(\Phi\rvert$ for
          some $J>0$ leaves the transfer matrix with one nonzero eigenvalue, equal
          to one, hence a single normal block filling the ambient bond space.
          This is the spectral step of
          \cite[lines~344--355]{Cirac2017MPU} run from the stabilized power
          instead of from the shifted traces. Unitarity of the periodic
          family is a premise and not a conclusion of the spectral step: the
          bundled presentation contains it, while an exact rank-one transfer
          power constrains only the normalized double layer. Every entry the
          converse is meant to serve carries it already.

          Trace normalization need not be added to the generic supplied-pair
          entry. Write $P=\lvert\rho)(\Phi\rvert$ and $t=\tr(\rho)$.
          Since $(\Phi\vert\rho)=t$, direct multiplication gives
          \begin{align}
              P^2 & =tP, & \tr(P) & =t.
              \label{eq:mpu_canonical_form_full_support_pair_trace}
          \end{align}
          The MPU trace identities hold at every length greater than one
          \cite[lines~349--354]{Cirac2017MPU}. Thus, if $E^K=P$ and $K\geq2$,
          \begin{align}
              t & =\tr(P)=\tr(E^K)=1.
              \label{eq:mpu_canonical_form_full_support_trace_large_exponent}
          \end{align}
          If $K=1$, use lengths two and three instead of assuming a
          length-one trace identity. By
          (\ref{eq:mpu_canonical_form_full_support_pair_trace}),
          \begin{align}
              t^2 & =\tr(P^2)=\tr(E^2)=1, &
              t^3 & =\tr(P^3)=\tr(E^3)=1, & t & =1.
              \label{eq:mpu_canonical_form_full_support_trace_exponent_one}
          \end{align}
          The exponent premise costs nothing. The generic supplied-pair entry
          admits the exponent zero, where the condition reads
          $\mathbf 1=\lvert\rho)(\Phi\rvert$; the right side has rank at most
          one. The bond dimension cannot vanish, since a zero-dimensional
          bond space gives a zero periodic operator on the nonempty physical
          space. Hence $D=1$, and taking traces gives
          \begin{align}
              t & =\tr(P)=\tr(\mathbf 1_{D^2})=D^2=1.
              \label{eq:mpu_canonical_form_full_support_trace_exponent_zero}
          \end{align}
          Thus the normalization required by the converse follows in every
          case. In bond dimension one the normalized transfer matrix of a
          tensor with a unitary periodic
          family is already the identity
          \cite[lines~397--409]{Cirac2017MPU}. The same pair then satisfies
          the condition at the exponent one, so the converse at a positive
          exponent covers the zero-exponent inputs as well.

          The positive-exponent constructor
          \leanid{MPOTensor.IsMPU.canonicalFormIIOfTransferPower} and its
          arbitrary-exponent extension
          \leanid{MPOTensor.IsMPU.canonicalFormIIOfSuppliedTransferPower}
          turn a supplied fixed pair into the bundled presentation of the
          tensor that carries it, and nothing more. So it removes the
          restriction exactly where the
          tensor carrying the datum is the tensor the standing convention
          already governs: the four entries
          \texttt{lem:mpu\_admissible\_source\_u\_isometry} and the three
          former candidates. It dissolves the citation obstruction above:
          the eleven retained entries obtain the presentation from their own
          supplied pair, with the same $\rho$, and the three duplicates are
          removed. The second item is not needed for them. For the twelve
          it rewrites the restriction without weakening it, because the
          tensor carrying the datum is a derived one and the converse
          alone constructs no datum for a derived tensor. Positive blocking,
          physical adjunction, and independent tensor products are supplied
          separately by preservation constructors, without identifying source
          factors across tensors.
    \item Preservation of the bundled presentation along each operation the
          twelve use: positive physical blocking, physical adjunction,
          transposition, complex conjugation, independent tensor products,
          and composition. Positive blocking is now proved by
          \leanid{MPOTensor.IsMPUCanonicalFormII.blockTensor}, and physical
          adjunction by
          \leanid{MPOTensor.IsMPUCanonicalFormII.physicalAdjointTensor}.
          Independent tensor products are covered by
          \leanid{MPOTensor.IsMPUCanonicalFormII.tensorProduct}. The other
          operations remain open. A preservation statement must deliver all four
          clauses of the presentation for the derived tensor: unitarity of
          its periodic operator family, literal canonical-form-II data for
          its normalized flattening, full support of those data, and a
          positive diagonal trace-one ambient matrix fixed by its normalized
          transfer map. Blocking is the case the source states outright at
          line~356. For composition no such statement about the product
          tensor is available, for the reason given below; that case has to
          go through a reduced representative instead.
\end{itemize}

The second item is not uniformly distant. The four clauses have been assembled
for positive blocking: unitarity, canonical-form-II data for the blocked
normalized flattening, full support, and preservation of the recorded
positive diagonal fixed matrix. In particular, whenever the blocked tensor is
simple, its source-$v$ gate is isometric for that same matrix by
\leanid{MPOTensor.IsMPUCanonicalFormII.sourceV_blockTensor_isIsometry}; a
simple block occurs at a positive length at most $D^4$ by
\leanid{MPOTensor.IsMPUCanonicalFormII.exists_sourceV_blockTensor_isIsometry}.
Neither result transports source-cut factors back to the unblocked tensor.
Physical adjunction is also assembled: it preserves unitarity, transports the
data together with their full support, and uses diagonality to identify the
conjugated ambient fixed matrix with the original one. Its bundled
preservation theorem is
\leanid{MPOTensor.IsMPUCanonicalFormII.physicalAdjointTensor}. For
independent tensor products, the four clauses are also assembled by
\leanid{MPOTensor.IsMPUCanonicalFormII.tensorProduct}. The ambient matrix is
the reindexed Kronecker product of the two recorded matrices. It is positive
definite, diagonal, of trace one, and fixed by the product transfer map.
These preservation results do not identify raw source cuts or chosen
compact-SVD factors across tensors. Transposition and conjugation are
transported only through their composite, physical adjunction; neither half
is transported on its own.

Composition differs from the other derived-tensor operations. Unitarity of a
composition is available, but canonical-form-II data need not exist on the
whole product bond space. The representative-index additivity theorem avoids
such a preservation claim: it reduces the composition tensor and compares
source ranks by inequalities, then uses the exact simple-tensor rank product
to force saturation. This does not identify selected source factors or raw
source cuts before and after reduction.

A composition keeps the whole product bond space $D_1D_2$, and that space can
be too large for the full-support clause. Take the right shift on $d\geq 2$
levels and its physical adjoint, both matrix product unitaries
(\leanid{MPOTensor.rightShiftTensor} and \leanid{MPOTensor.leftShiftTensor}).
Their composite generates the identity operator at every length, and in the
product bond space $\mathbb C^{d}\otimes\mathbb C^{d}$ its one-site matrices
are
\begin{align}
    A^{(i,k)} & =\lvert i,k\rangle\langle\Phi\rvert,\qquad
    \langle\Phi\rvert=\sum_{\gamma}\langle\gamma,\gamma\rvert.
    \label{eq:mpu_canonical_form_full_support_shift_composite}
\end{align}
Each of these has rank one, and they share the row functional
$\langle\Phi\rvert$, so every product of them, at every length, again has rank
one and every element of the span of the length-$L$ products has rank at most
one. A tensor with full-support canonical-form-II data is a unitary conjugate
of a direct sum of nonzero multiples of normal blocks whose dimensions add up
to the ambient dimension. A retained block of dimension at least two would put
a matrix of rank at least two into that span, so every block would have
dimension one; but then the one-site matrices would all be diagonal in one
basis and would commute, and
Eq.~\ref{eq:mpu_canonical_form_full_support_shift_composite} gives
$A^{(0,0)}A^{(0,1)}=0$ while
$A^{(0,1)}A^{(0,0)}=\lvert 0,1\rangle\langle\Phi\rvert\neq 0$.
The composite therefore has no
full-support canonical-form-II data in the product bond space at all, and no
amount of blocking repairs it, since the blocked matrices are the same
products and stay of rank one.

This is the enlargement phenomenon of the opening section met again: the
composite is a nonminimal ambient representative of the operator family it
generates, and the presentation is a property of the representative and not of
the operators. The composition case therefore does not ask for a preservation
statement about the product tensor. The public index theorem applies
\leanid{MPOTensor.IsMPU.exists_reduced_cfii_representative} to the product.
Equality of positive-length periodic operators supplies a rectangular
reduction from the product to that representative
(\leanid{MPOTensor.exists_isReduction_of_mpo_eq_cfii}). The inequalities
\leanid{MPOTensor.rightRank_blockTensor_le_of_isReduction} and
\leanid{MPOTensor.leftRank_blockTensor_le_of_isReduction} carry the
four-site composition bounds to the representative. Because its simple
four-site block has the maximal possible rank product, both bounds are
saturated (\leanid{MPOTensor.IsMPUCanonicalFormII.compositionRanks_of_isReduction}).
This proves the representative-index additivity without identifying raw
source cuts across reduction. The printed composition argument blocks its
composite until it is simple \cite[lines~837--841]{Cirac2017MPU}; for a
nonminimal product presentation such as the example above, blocking alone
cannot supply that simple tensor, and a reduction is required. The old
six-block assumptions are not hypotheses of the proved composition entry.
They still appear in the unconsolidated admissible Index Theorem, but the
composition step no longer needs them.

The first item and the three-label consolidation are complete. Thirteen
pointwise entries remain: the generic supplied-pair step, which is kept
deliberately, and the twelve derived-tensor entries. Their positive-blocking,
physical-adjoint, and independent-product cases now have bundled preservation
theorems, but the downstream entries have not yet been consolidated. Separate
conjugation and transposition comparisons, symmetry transforms, and claims
about selected factors of a composition retain distinct construction
obligations; representative-index composition additivity does not.
The twenty-three pathwise entries retain their continuous-source-data hypotheses.

\section{Verdict}

\textbf{Verdict: scope restriction.} The normal-tensor proposition,
source-isometry lemma, and Theorem~III.8 of Ref.~\cite{Cirac2017MPU} concern
the intended reduced canonical-form-II representative with the ambient zero
complement removed. The normal-tensor proposition does not state this reduction
as a hypothesis, and literal canonical-form membership alone does not imply it.

A source-faithful formalization must therefore construct the reduced
representative. If a theorem retains named source cuts, the formalization must
also justify the corresponding source constructions separately. The
representative reduction is now formalized for every positive bond dimension
and preserves the positive-length periodic operator family. Statements about a
tensor already in canonical form II take the bundled presentation, which is the
source's own standing hypothesis. Statements that would start from an arbitrary
ambient MPU still need an identification of source factors across potentially
unequal bond dimensions, which has not been proved.

Within Theorem~III.8 itself, the complete source-$u$ contraction and the
complete source-$v$ equivalence are formalized. Both are stated for a tensor
satisfying the bundled canonical-form-II presentation, which is exactly the
hypothesis the source states once before Lemma~III.7 and never repeats. The
ambient positive diagonal weight and the stabilized rank-one transfer power
$E^{J}=\lvert\rho)(\Phi\rvert$ at the positive exponent $J=\max(D^2-1,1)$ are
consequences of that presentation rather than supplied data, so
\texttt{lemuisometry} and \texttt{ThmFund1} carry the source's hypothesis set.
The generic supplied-fixed-pair contractions remain available as the reusable
step, stated for an arbitrary positive diagonal weight and an arbitrary
stabilization exponent.

What remains open is the transport of source cuts. For an arbitrary MPU
representative the reduction produces a new tensor of possibly smaller bond
dimension, and no theorem identifies the compact-SVD cuts, ranks, or selected
source factors of that representative with those of the original ambient
tensor. That transport remains in \ghissue{7012}, and the unrestricted
representative remains in \ghissue{7127}. Downstream standard-form, the
unreduced source-cut index formula, continuity, and symmetry results also lack
existence or preservation theorems for selected source data under products,
transformed comparison tensors, and continuous paths. The representative
index and its composition additivity do not depend on those transports.
Chapter~28 records these requirements as explicit admissibility hypotheses
rather than consequences of the MPU predicate. They come off in two stages,
neither singly nor all at once. The converse of the stabilization theorem is
proved, and the three same-tensor duplicates have been consolidated with
their formerly obstructing citations. Preservation along each remaining
derived-tensor operation releases the restrictions on a physical block, a
comparison tensor, or a tensor product. Composition additivity of the
representative index instead uses a reduced representative, rectangular rank
inequalities, and rank saturation. The elimination plan above records the
separate tasks.

The full-support irreducibility implication is formalized by
\leanid{MPSTensor.CPSVCanonicalFormData.isIrreducibleTensor_of_r_eq_one_of_fullSupport}.
The shifted normalized transfer traces give the spectral-radius and primitivity
conclusions, which are combined with full support in
\leanid{MPOTensor.IsMPU.isNormalTensor_normalizedFlattening_of_cpsv}.
The formalization deliberately does not use bare
\leanid{MPOTensor.IsMPU} to infer normality or an ambient positive-definite
fixed point for the original, potentially nonminimal representative. Instead,
it uses
\leanid{MPOTensor.IsMPU.exists_reduced_cfii_representative} to construct the
reduced representative.

\bibliographystyle{alpha}
\bibliography{references}
\end{document}
